%\section{How is cointegration tested?}
%\section{Cointegration test}
\section{How?}

\begin{frame}[allowframebreaks]
	\frametitle{Johansen test}
	\small
	\begin{itemize}
		\item \textbf{Goal}: We need a formal test to find out if cointegration exists and, if so, how many (i.e., how many ``leashes").
		
		\item \textbf{Test}: The Johansen test (Section 5.9) is the standard tool.
		
		\item \textbf{The model}: Vector Error-Correction Model (VECM) form (from Section 5.6):
		$$\Delta z_t = \boldsymbol{\Pi} z_{t-1} + \sum_{i=1}^{p-1} \boldsymbol{\Phi_i}^* \Delta z_{t-i} + c(t) + \boldsymbol{\Theta(B)}a_t$$
		
		\item \textbf{Objective}: To check for the rank ($m$) of the $\boldsymbol{\Pi}$ matrix which is the number of cointegrating vectors.
		\begin{itemize}
			\item If $\text{Rank}(\boldsymbol{\Pi}) = 0$, the series are not cointegrated.
			
			\item If $\text{Rank}(\boldsymbol{\Pi}) = m > 0$, there are $m$ cointegrating vectors.
		\end{itemize}
	\end{itemize}
	\framebreak
	
	\begin{itemize}
		\item \textbf{Test statistics}: The test produces two statistics (Section 5.9.4):
		\begin{enumerate}
			\item Trace statistic ($L_{tr}$): Tests 
			$$H_0: \text{Rank} = m\ vs.\ H_a: \text{Rank} > m$$
			
			\item Max-Eigenvalue statistic ($L_{max}$): Tests 
			$$H_0: \text{Rank} = m\ vs.\ H_a: \text{Rank} = m+1$$
		\end{enumerate}
		
		\item $L_{tr}(m) = -(T - kp) \sum_{i=m+1}^{k} \ln(1 - \lambda_i)$
		
		\item $L_{max}(m) = -(T - kp) \ln(1 - \lambda_{m+1})$
		\item We keep testing sequentially for all values of \textit{m} until it reaches Rank($\boldsymbol{\Pi}$).
	\end{itemize}
\end{frame}

\begin{frame}[allowframebreaks]
	\frametitle{Example}
	\small
		
	\textbf{Setup:}
	We have two nonstationary (\(I(1)\)) series \(z_t = (z_{1t}, z_{2t})'\), and we wish to test for the number of cointegrating relationships.
	
	\textbf{Step 1: Set up the ECM Model}
	Choose a VAR order \(p = 2\). Based on the text, the Error-Correction Model (ECM) is:
	\[
	\Delta z_t = c_0 + \boldsymbol{\Pi} z_{t-1} + \boldsymbol{\Phi}_1^* \Delta z_{t-1} + a_t
	\]
	Our goal is to determine the \emph{rank} (\(m\)) of the \(2\times2\) matrix \(\boldsymbol{\Pi}\):
	\begin{itemize}
		\item \(m = 0\): no cointegration,
		\item \(m = 1\): one cointegrating vector,
		\item \(m = 2\): both series are stationary.
	\end{itemize}
	
	\framebreak
	\textbf{Step 2: Run ``Helper'' Regressions}
	We cannot test \(\boldsymbol{\Pi}\) directly. We first concentrate out short-run dynamics:
	\begin{enumerate}
		\item Regress \(\Delta z_t\) on a constant and \(\Delta z_{t-1}\); save residuals \(\hat{u}_t\).
		\item Regress \(z_{t-1}\) on a constant and \(\Delta z_{t-1}\); save residuals \(\hat{v}_t\).
	\end{enumerate}

	\framebreak
	\textbf{Step 3: Compute Covariance Matrices}
	From the residuals, we compute sample covariance matrices (hypothetical values):
	\[
	\hat{\Sigma}_{00} = 
	\begin{bmatrix}
		4 & 1\\
		1 & 3
	\end{bmatrix}, \quad
	\hat{\Sigma}_{11} =
	\begin{bmatrix}
		5 & 2\\
		2 & 2
	\end{bmatrix}, \quad
	\hat{\Sigma}_{01} =
	\begin{bmatrix}
		3 & 1.5\\
		1 & 1
	\end{bmatrix}, \quad
	\hat{\Sigma}_{10} = \hat{\Sigma}_{01}' =
	\begin{bmatrix}
		3 & 1\\
		1.5 & 1
	\end{bmatrix}
	\]
	
	\framebreak
	\textbf{Step 4: Solve the Eigenvalue Problem}
	The Johansen test uses the eigenvalues of
	\[
	S = \hat{\Sigma}_{11}^{-1}\hat{\Sigma}_{10}\hat{\Sigma}_{00}^{-1}\hat{\Sigma}_{01}.
	\]
	
	\framebreak
	\textbf{Calculations:}
	\begin{enumerate}
		\item Find inverses:
		\[
		|\hat{\Sigma}_{00}| = 11, \quad 
		\hat{\Sigma}_{00}^{-1} = \tfrac{1}{11}
		\begin{bmatrix}
			3 & -1\\
			-1 & 4
		\end{bmatrix}, \quad
		|\hat{\Sigma}_{11}| = 6, \quad
		\hat{\Sigma}_{11}^{-1} = \tfrac{1}{6}
		\begin{bmatrix}
			2 & -2\\
			-2 & 5
		\end{bmatrix}
		\]
		\item Multiply to obtain:
		\[
		S =
		\begin{bmatrix}
			0.3636 & 0.1591\\
			0.2273 & 0.1932
		\end{bmatrix}
		\]
		\item Eigenvalues of \(S\):
		\[
		\lambda_1 \approx 0.487, \qquad \lambda_2 \approx 0.070
		\]
	\end{enumerate}
	
	\framebreak
	\textbf{Step 5: Compute Test Statistics}
	
	Effective sample size: \(T - kp = 196\).
	
	\textbf{Trace Statistic:}
	$$L_{\mathrm{tr}}(0) = -196 \left[\ln(1 - 0.487) + \ln(1 - 0.070)\right]$$
	$$= -196[-0.667 - 0.072] \approx \mathbf{144.8}$$
	
	$$L_{\mathrm{tr}}(1) = -196 \ln(1 - 0.070)
	= -196[-0.072] \approx \mathbf{14.1}$$
	
	\textbf{Max-Eigenvalue Statistic:}
	$$L_{\max}(0) = -196 \ln(1 - 0.487)
		= -196[-0.667] \approx \mathbf{130.7}$$
		
	$$L_{\max}(1) = -196 \ln(1 - 0.070)
	= -196[-0.072] \approx \mathbf{14.1}$$
	
	\framebreak
	\textbf{Step 6: Compare with Critical Values}
	\centering
	\renewcommand{\arraystretch}{1.2}
	\begin{tabular}{lcccc}
		\hline
		\textbf{Hypothesis} & \textbf{Statistic} & \textbf{Value} & \textbf{5\% CV} & \textbf{Decision}\\
		\hline
		\multicolumn{5}{l}{\textit{Trace Test}}\\
		$H_0: m=0$ & $L_{\mathrm{tr}}(0)$ & 144.8 & 17.95 & Reject $H_0$\\
		$H_0: m \le 1$ & $L_{\mathrm{tr}}(1)$ & 14.1 & 8.18 & Reject $H_0$\\[2mm]
		\multicolumn{5}{l}{\textit{Max-Eigen Test}}\\
		$H_0: m=0$ & $L_{\max}(0)$ & 130.7 & 14.90 & Reject $H_0$\\
		$H_0: m=1$ & $L_{\max}(1)$ & 14.1 & 8.18 & Reject $H_0$\\
		\hline
	\end{tabular}
	
	\framebreak
	\textbf{Conclusion}
	Since both tests reject \(H_0\) for \(m=0\) and \(m=1\), we conclude that the rank is \(m=2\).  
	Therefore, the original bivariate system \(z_t\) is stationary (\(I(0)\)) and has full rank.
\end{frame}